\documentclass[10pt,a4paper]{amsart}
\usepackage[margin=19mm]{geometry}
\usepackage{amsmath,amssymb,mathtools}
\usepackage[scr=rsfs,cal=euler]{mathalfa}
\usepackage{xfrac}
\usepackage[unicode,hidelinks,bookmarks=false]{hyperref}
\newcommand{\ie}{i.e.\ }
\newcommand{\cf}{cf.\ }
\newcommand{\resp}{respectively\ }
\newcommand{\Subsection}{\subsection}
\providecommand{\nobreakdash}{\nobreak\hbox{-}}
\setlength{\emergencystretch}{2em}
\multlinegap0pt
\usepackage{bm}
\usepackage{bbm}
\usepackage{dsfont}
\usepackage{xspace}
\usepackage{cancel}
\usepackage{mathtools}
\usepackage{amsthm}
\makeatletter
\@ifundefined{thm}{\newtheorem{thm}{Thm}}{}
\@ifundefined{prop}{\newtheorem{prop}[thm]{Proposition}}{}
\makeatother
\makeatletter
\def\ch{\mathcal{H}}
\expandafter\ifx\csname thmA\endcsname\relax
\newenvironment{thmA}[1]{%
  \renewcommand\thenamedthm{#1}%
  \begin{namedthm}
}{%
  \end{namedthm}
}
\fi
\makeatother
\title[Centralizers in Hecke Algebras of Any Coxeter Group]{Centralizers in Hecke Algebras of Any Coxeter Group}
\author{Haiyu Chen}
\date{Source: arXiv:2508.14473v1 (CC BY 4.0)}
\begin{document}
\maketitle

\noindent\emph{Source and licence.} Centralizers in Hecke Algebras of Any Coxeter Group. Version arXiv:2508.14473v1 (CC BY 4.0), distributed under the Creative Commons Attribution 4.0 International licence, as recorded in the arXiv licence metadata for this exact version and shown on the abstract page https://arxiv.org/abs/2508.14473v1. Full rights notice: licenses/arxiv-2508.14473-CC-BY-4.0.txt.\par\smallskip

\noindent\textit{Editorial scope.} Counted unit: the Prop whose source label is \texttt{prop: diamond} in the article (source0.tex lines 681--687, source label \texttt{prop: diamond}), delivered together with the article's own complete original proof of it (source0.tex lines 689--706, one \texttt{proof} environment of 18 lines). No mathematics is written, repaired or re-derived: statement and proof are the article's own text line for line. Required context retained uncounted: none. Every definition, display and label of those lines is retained unchanged. Omitted from this package and not redistributed: 1--680, 688--688, 707--1140. Nothing inside a retained excerpt is omitted or altered: every retained body is delivered exactly as the article's lines are written, so no omission receipt of any kind accompanies this package. Citations: the retained excerpts carry no citation command at all, so no bibliography entry of the article is transcribed and no reference is redistributed. Reference adaptation: none was needed and none was made; every reference inside the delivered unit resolves inside it, so no unresolved reference or citation is produced. Printed slips kept literal: no printed word, symbol, hypothesis or number of the counted statement or of its proof was altered. Layout and class: the article's own class is \texttt{amsart} with packages including color, amsthm, amsfonts, amssymb, amsmath, amsxtra, xy, xr-hyper, hyperref, verbatim, enumerate, cleveref, geometry, mathrsfs; the wrapper is the lane's pinned \texttt{amsart} editorial wrapper at 10pt on A4, which restates the theorem-like environments the retained text needs and adds the article's own macro definitions that the retained text needs (\texttt{\textbackslash ch}), with extra wrapper packages (bm, bbm, dsfont, xspace, cancel, mathtools). The delivered numbering is the unit's own. Assets: no retained span contains a figure environment, a graphics-inclusion command, a PDF-page inclusion, a table or a TikZ picture; no image file is copied into this package, the package declares no asset dependency and no image file is redistributed with it. Licence: version arXiv:2508.14473v1 of this article is distributed under the Creative Commons Attribution 4.0 International licence, as recorded in the arXiv licence metadata for this exact version and shown on the abstract page https://arxiv.org/abs/2508.14473v1. A full rights notice is delivered at licenses/arxiv-2508.14473-CC-BY-4.0.txt. Characters: every retained excerpt is pure ASCII, so the delivered engine, XeLaTeX with its default Latin Modern text font, reports no missing character.
        \begin{prop}\label{prop: diamond}
            Write $x=\sum_{w\in W}x_wT_w$, $x_w\in R$. Then $x\in \mathcal{Z}_{\ch}(\ch_{J})$ if and only if
            \begin{enumerate}
                \item $x_w=x_{w'}$ whenever $w\approx_J w'$;
                \item $x_{w'}=b_sx_w-a_sx_{sw}$ whenever $w\xrightarrow{s} w'$ and $\ell(w)>\ell(w')$ for $s\in J$.
            \end{enumerate}
        \end{prop}

        \begin{proof}
            For $z\in \ch$ and $w\in W$, we use $(z)_w$ to denote the coefficient of the basis element $T_w$ in $z$, for example, $(x)_w=x_w$.

            We first show the only if part.
            \begin{enumerate}
                \item It suffices to show $x_w=x_{w'}$ for $w\xrightarrow{s} w'$ with $\ell(w')=\ell(w)$. We may assume $w'\neq w$.
                
                By the exchange property of Coxeter groups, we cannot have both $\ell(sw)<\ell(w)$ and $\ell(ws)<\ell(w)$. Indeed, $\ell(sw)<\ell(w)$ implies there must exist a reduced expression $w$ starting with $s$, say $w=s_0s_1s_2\cdots s_k$ with $s=s_0$. But then $\ell(ws)<\ell(w)$ implies $ws=s_0s_1\cdots \hat{s_j}\cdots s_k$ with $\hat{}$ denoting omission. If $j\neq 0$ then $\ell(sws)<\ell(w)$. If $j=0$, then $ws=sw$ and that further implies $w'=sws=ssw=w$.
                
                Hence either $\ell(sw)>\ell(w)=\ell(w')$ or $\ell(ws)>\ell(w)=\ell(w')$.

                In the first case, $(T_sx)_{sw}=x_w+a_sx_{sw}$ and $(xT_s)_{sw}=x_{sws}+a_sx_{sw}$. Hence $x_{w'}=x_w$. The second case is similar.

                \item In this case, we have $\ell(sws)<\ell(ws)=\ell(sw)<\ell(w)$ and $sw\approx_J ws$. Then $(T_sx)_{sw}=b_sx_w$ and $(xT_s)_{sw}=a_sx_{sw}+x_{sws}$. Hence $x_{w'}=b_sx_w-a_sx_{sw}$. 
            \end{enumerate}

            Now we show the if part, i.e., for $s\in J$, we want to show $T_sx=xT_s$ when $x_w$, $w\in W$ satisfy both conditions. For each $w\in W$, $s\in J$, there are two cases. If $|\ell(sws)-\ell(w)|=2$ or ($\ell(sws)=\ell(w)$ and $|\ell(ws)-\ell(sw)=2|$), we are in the situation of $(ii)$. By the proof in $(i)$ above, the case $\ell(w)=\ell(sws)>\ell(ws)=\ell(sw)$ is trivial. The only case left is $\ell(w)=\ell(sws)<\ell(ws)=\ell(sw)$. But then this implies $sws\approx_J w$, and $sw\approx_J ws$. We are in the situation of $(i)$. Hence $T_sx=xT_s$ for any $s\in J$.
        \end{proof}

\end{document}
